\section{Motivating Traces}
\label{sec:motivating}

The three traces below are not hypotheticals. Each is a run of our harness
(\cref{sec:evaluation}), replayed from its event log and cross-checked against
the mock provider's ground-truth ledger --- an independently written record of
every mutation the provider \emph{actually applied}. In each, the agent
workload is the same: acquire work, call one non-idempotent endpoint, record
the outcome.

\subsection{Trace 1 --- the duplicate nobody sees (B0, naive retry)}

A worker sends the mutation. The provider applies it. The worker is
\texttt{SIGKILL}ed after the response is on the wire but before the outcome is
recorded (crash point \texttt{after\_response\_before\_resolution}). A
supervisor observes an execution with no terminal record and runs the step
again. The provider applies the mutation a second time.

The ground-truth ledger now holds two applied mutations for one logical
step. The system's own records hold one successful execution. Nothing in the
system is in a state that a monitor could alert on: there is no error, no
retry counter above its threshold, no ambiguous record. The execution is
\emph{green}.

This is not a strawman configuration; it is what a retry-on-timeout loop does,
and it is the behaviour we measure for B0. Across the crashed cells B0 ends
this way in \BaselineDupLow{}--\BaselineDupHigh{} of executions depending on
the endpoint, and so do B1 and B2 (\cref{tab:outcomes}).
% \BaselineDupLow / \BaselineDupHigh: generated/numbers.tex, the min and max
% undetected-duplicate rate over {B0,B1,B2} x three response classes,
% crashed regime.

\subsection{Trace 2 --- a lease and a fenced write do not help}

The same crash, against B1 (a real Redis lease) and B2 (a lease plus an exact
expected-version compare-and-swap on the state document). Both are the patterns
an experienced engineer reaches for, and both are doing their job: no stale
writer corrupts the state document, and no two workers hold the lock at once.

Neither prevents the second effect. The lease governs \emph{who may write
state}; the CAS governs \emph{which version of state may be replaced}. The
external call is neither. When the supervisor re-runs the step it holds a
valid, freshly acquired lease and writes a correctly versioned record --- of a
duplicate. \Cref{sec:evaluation} shows B1 and B2 duplicating at
substantially the same rate as B0, which is the point: the mechanisms that
protect the database do not protect the world.

\subsection{Trace 3 --- a durable log is necessary and not sufficient (B4)}

B4 is an event-sourced durable-execution engine of the kind these systems are
built on: an append-only history in Redis, each append acknowledged durable
with the \emph{same} \texttt{WAITAOF} barrier AEP uses, replay after a crash,
and completed activities memoised so they are not re-run. It is not ablated on
durability. Its history is on disk before the call.

It duplicates anyway. The history records that the activity was
\emph{scheduled}; it does not, and cannot, record whether the provider applied
the effect, because that fact was never available to the engine. On replay the
engine sees a scheduled activity with no completion and --- following the
vendor's documented default, where Activity Maximum Attempts is
unlimited~\cite{temporal-retry} --- schedules it again.

The mitigation the vendor recommends is to make the activity
idempotent~\cite{temporal-activity-definition}. That is exactly the
precondition this paper's premise removes.

\subsection{Trace 3b --- and the at-most-once configuration loses the effect}

Setting Maximum Attempts to 1 is documented and standard: ``Setting the value
to 1 means a single execution attempt and no retries''~\cite{temporal-retry}.
We run it as B4b. It cannot produce a caller-caused duplicate, and it does not.

What it produces instead is the other silent failure. Where the provider
applied the mutation before the worker died, B4b's history records an activity
that timed out; the workflow fails; the effect exists and the records deny
it. Critically, that record is an ordinary failure and not a declared
ambiguity --- there is no outcome class in B4b in which the engine can say
\emph{I do not know}, so nothing escalates and no operator is told that a
real-world effect may be unaccounted for.

\subsection{What the traces establish}

\begin{table}[t]
\centering
\caption{The trilemma, as the systems under comparison instantiate it. A
durable write-ahead record is necessary (Traces 1, 2) and not sufficient
(Trace 3); what distinguishes AEP is an outcome class in which it can decline
to decide. B3 --- AEP with the durability barrier removed and nothing else ---
reaches the same corner, which is the point: the barrier is not what buys
declared ambiguity (\cref{sec:eval-detection} measures the ablation), and
\cref{sec:eval-prevention} isolates what it does buy.}
\label{tab:trilemma}
\small
\begin{tabular}{@{}lccc@{}}
\toprule
 & \undetdup{} & \lostfx{} & \declamb{} \\
\midrule
B0 naive retry            & \textbf{yes} & --- & no \\
B1 lease-only             & \textbf{yes} & --- & no \\
B2 CAS-only               & \textbf{yes} & --- & no \\
B4 durable, max attempts $\infty$  & \textbf{yes} & --- & no \\
B4b durable, max attempts $=1$     & no & \textbf{yes} & no \\
\midrule
B3 AEP minus the barrier  & no & no & \textbf{yes, bounded} \\
AEP-full                  & no & no & \textbf{yes, bounded} \\
\bottomrule
\end{tabular}
\end{table}

\Cref{tab:trilemma} is the shape of the argument. At-most-once dispatch is
available off the shelf --- B4b buys it with one configuration setting, and so
does AEP's ablation B3. What is not available off the shelf is knowing
\emph{which} of duplicate-or-loss occurred, because the fact a system would
need for that --- whether the provider applied the effect --- is not in any log
it can write and cannot be put there by any amount of logging. It has to be
obtained from the endpoint, and \cref{sec:eval-rq1} shows the price of that: the
achievable outcome is bounded by what the endpoint can be asked.
